\documentclass[11pt,a4paper]{article}
\usepackage[utf8]{inputenc}
\usepackage[T1]{fontenc}
\usepackage[ngerman]{babel}
\usepackage{hyperref}
\usepackage{geometry}
\usepackage{amsmath}
\usepackage{amssymb}
\usepackage{listings}
\usepackage{xcolor}
\usepackage{booktabs}
\usepackage{enumitem}
\usepackage{graphicx}
\usepackage{fancyhdr}
\usepackage{titlesec}
\usepackage{longtable}
\usepackage{array}
\usepackage{mdframed}
\usepackage{multirow}

\geometry{margin=2.5cm, headheight=14pt}
\setcounter{tocdepth}{3}

% -----------------------------------------------------------------------
% Code-Listing-Stil
% -----------------------------------------------------------------------
\definecolor{codebg}{rgb}{0.95,0.95,0.97}
\definecolor{codegreen}{rgb}{0.0,0.5,0.0}
\definecolor{codegray}{rgb}{0.4,0.4,0.4}
\definecolor{codeblue}{rgb}{0.0,0.0,0.7}
\definecolor{codepurple}{rgb}{0.5,0.0,0.5}
\definecolor{noteboxbg}{rgb}{0.95,0.98,1.0}
\definecolor{noteboxborder}{rgb}{0.2,0.5,0.8}
\definecolor{warnboxbg}{rgb}{1.0,0.97,0.90}
\definecolor{warnboxborder}{rgb}{0.8,0.5,0.0}

\lstdefinestyle{python}{
  language=Python,
  backgroundcolor=\color{codebg},
  basicstyle=\ttfamily\small,
  keywordstyle=\color{codeblue}\bfseries,
  stringstyle=\color{codegreen},
  commentstyle=\color{codegray}\itshape,
  numberstyle=\tiny\color{codegray},
  numbers=left,
  numbersep=8pt,
  frame=single,
  rulecolor=\color{codegray!40},
  breaklines=true,
  showstringspaces=false,
  tabsize=4,
  captionpos=b,
  morekeywords={self, True, False, None, async, await},
}

\lstdefinestyle{batch}{
  language={},
  backgroundcolor=\color{codebg},
  basicstyle=\ttfamily\small,
  commentstyle=\color{codegray}\itshape,
  numbers=left,
  numberstyle=\tiny\color{codegray},
  numbersep=8pt,
  frame=single,
  rulecolor=\color{codegray!40},
  breaklines=true,
  captionpos=b,
}

\lstset{style=python}

\newmdenv[
  backgroundcolor=noteboxbg,
  linecolor=noteboxborder,
  linewidth=1.5pt,
  leftline=true, rightline=false, topline=false, bottomline=false,
  innerleftmargin=8pt,
  skipabove=6pt, skipbelow=6pt
]{notebox}

\newmdenv[
  backgroundcolor=warnboxbg,
  linecolor=warnboxborder,
  linewidth=1.5pt,
  leftline=true, rightline=false, topline=false, bottomline=false,
  innerleftmargin=8pt,
  skipabove=6pt, skipbelow=6pt
]{warnbox}

% -----------------------------------------------------------------------
% Kopf- und Fußzeile
% -----------------------------------------------------------------------
\pagestyle{fancy}
\fancyhf{}
\rhead{Dobot ROS2 Bridge -- Isaac Sim Extension}
\lhead{Dobot ROS2 Bridge}
\rfoot{Seite \thepage}

% -----------------------------------------------------------------------
% Dokument
% -----------------------------------------------------------------------
\title{
  \textbf{Dobot ROS2 Bridge}\\[6pt]
  \large Isaac Sim Extension -- Technische Dokumentation\\[4pt]
  \normalsize Bidirektionaler Digital-Twin-Kopplungsrahmen für den\\
  Dobot Magician in NVIDIA Isaac Sim 5.1
}
\author{Tobias Küster}
\date{\today}

\begin{document}
\maketitle
\tableofcontents
\newpage

% ===================================================================
\section{Übersicht und Zielsetzung}
% ===================================================================

Die Extension \texttt{dobot\_ros} implementiert eine \textbf{bidirektionale
Echtzeit-Kopplung} zwischen einem realen \textbf{Dobot Magician}-Roboterarm
und einem \textbf{NVIDIA Isaac Sim 5.1}-Simulationsmodell. Sie stellt damit
einen vollständigen Digital-Twin-Kopplungsrahmen dar, der es ermöglicht,
den physischen Roboter in einer virtuellen Simulationsumgebung zu spiegeln
und umgekehrt Steuerbefehle aus der Simulation an die reale Hardware zu
senden.

\subsection{Zielsetzung}
\begin{itemize}
  \item \textbf{Digital Twin (Richtung 1 -- Real $\to$ Sim):}
        Der echte Roboter bewegt sich, und das Simulationsmodell folgt
        in Echtzeit. Gelenkwinkel werden per USB über die
        \texttt{pydobot}-Bibliothek gelesen, kinematisch transformiert
        und über die Isaac-Sim-Dynamic-Control-API auf die
        Simulationsgelenkstruktur übertragen.
  \item \textbf{Fernsteuerung (Richtung 2 -- Sim $\to$ Real):}
        Über das UI-Panel werden Jog-Befehle, Vertikalbewegungen,
        Vakuumgreifer-Schaltung, RC-Servo-Ansteuerung sowie
        automatisierte Pick-\&-Place-Sequenzen an den echten
        Roboter gesendet.
  \item \textbf{ROS2-Integration:}
        Die Extension startet intern einen ROS2-Node in Isaac Sim und
        publiziert Gelenkzustände auf dem Topic \texttt{/joint\_states}
        (Typ: \texttt{sensor\_msgs/msg/JointState}). Externe ROS2-Knoten
        (z.\,B. MoveIt!) können dieses Topic abonnieren und die aktuelle
        Kinematik in Echtzeit empfangen.
  \item \textbf{Interaktives Teach-in:}
        Bewegungsabläufe werden durch manuelles Verfahren des
        Roboters aufgezeichnet und anschließend vollautomatisch mit
        Sicherheitslogik wiedergegeben.
  \item \textbf{Würfel-Tracking (Ziel-Tracking):}
        Ein interaktiv in der Isaac-Sim-Szene verschiebbarer
        Ziel-Würfel dient als Bewegungssollwert für den echten
        Roboterarm.
\end{itemize}

% ===================================================================
\section{Systemumgebung und Softwarestapel}
\label{sec:system}
% ===================================================================

Dieser Abschnitt dokumentiert vollständig die Hardware- und
Softwareumgebung, in der die Extension entwickelt und betrieben wird.

\subsection{Hardware}

\begin{center}
\begin{tabular}{ll}
  \toprule
  Komponente & Spezifikation \\
  \midrule
  CPU & Intel Core i7-13620H (13. Generation, Hybrid-Architektur) \\
  RAM & 32\,GB DDR5 \\
  GPU & NVIDIA GeForce RTX 4070 Laptop GPU, 8\,GB VRAM (WDDM) \\
  Betriebssystem & Windows 11 Home, Build 26200 \\
  Roboter & Dobot Magician, USB-Verbindung über COM3 \\
  \bottomrule
\end{tabular}
\end{center}

\subsection{NVIDIA-Treiber und CUDA}

\begin{center}
\begin{tabular}{ll}
  \toprule
  Software & Version \\
  \midrule
  NVIDIA GeForce Treiber (Game Ready) & 592.27 (2. März 2026) \\
  CUDA Version (maximale Unterstützung durch Treiber) & 13.1 \\
  CUDA Toolkit (installiert) & 12.1.66 \\
  \quad \texttt{nvcc} (CUDA Compiler Driver) & 12.1, Build 32415258 \\
  CUDA-Pfad & \texttt{C:\textbackslash{}Program Files\textbackslash{}NVIDIA GPU Computing Toolkit\textbackslash{}CUDA\textbackslash{}v12.1} \\
  Nsight Compute & 2023.1.0 \\
  Nsight Systems & 2022.4.2 / 2023.1.2 \\
  NVIDIA Nsight Visual Studio Edition & 25.4.0.25287 \\
  \bottomrule
\end{tabular}
\end{center}

\begin{notebox}
\textbf{Treiber-Mindestanforderung von Isaac Sim 5.1:}
Laut der mitgelieferten Datei \texttt{kit/driver-requirements.json}
werden NVIDIA-Treiber unterhalb von Version \textbf{537.58} (Windows)
vom RTX-Renderer blockiert. Der installierte Treiber 592.27 liegt
deutlich darüber und ist vollständig kompatibel.
\end{notebox}

\subsection{NVIDIA Isaac Sim}

\begin{center}
\begin{tabular}{ll}
  \toprule
  Parameter & Wert \\
  \midrule
  Isaac Sim Version & 5.1.0-rc.19+release.26219.9c81211b.gl \\
  Installationspfad & \texttt{C:\textbackslash{}NVIDIA\_Isaac\_Sim} \\
  Interne Python-Version & 3.11.13 (gebündeltes Kit-Python) \\
  Python-Executable & \texttt{C:\textbackslash{}NVIDIA\_Isaac\_Sim\textbackslash{}python.bat} \\
  Kit-Executable & \texttt{C:\textbackslash{}NVIDIA\_Isaac\_Sim\textbackslash{}kit\textbackslash{}kit.exe} \\
  \bottomrule
\end{tabular}
\end{center}

\begin{notebox}
Isaac Sim bringt eine \textbf{eigene, isolierte Python-3.11.13-Laufzeitumgebung}
mit. Diese ist vollständig von einer etwaigen System-Python-Installation
getrennt. Alle Extension-Skripte laufen in dieser gebündelten
Python-Umgebung -- \textbf{nicht} in der System-Python-Installation
(hier: Python 3.14.3 im Projekt-\texttt{.venv}).
\end{notebox}

\subsection{ROS2-Integration in Isaac Sim}

\textbf{ROS2 wird nicht separat installiert.} Isaac Sim 5.1 enthält eine
eingebettete ROS2-Bridge-Extension unter:
\begin{center}
  \texttt{C:\textbackslash{}NVIDIA\_Isaac\_Sim\textbackslash{}exts\textbackslash{}isaacsim.ros2.bridge}
\end{center}

Diese Bridge stellt ROS2 Humble \textbf{und} ROS2 Jazzy als
vorkompilierte Shared Libraries bereit:

\begin{center}
\begin{tabular}{ll}
  \toprule
  Parameter & Wert \\
  \midrule
  Unterstützte ROS2-Distributionen & Humble (Standard) / Jazzy \\
  RMW-Implementierung (Middleware) & \texttt{rmw\_fastrtps\_cpp} (Fast DDS) \\
  Standard-Distro-Umgebungsvariable & \texttt{ROS\_DISTRO=humble} \\
  Library-Pfad (Humble) & \texttt{isaacsim.ros2.bridge\textbackslash{}humble\textbackslash{}lib} \\
  rclpy-Pfad & \texttt{isaacsim.ros2.bridge\textbackslash{}rclpy} \\
  \bottomrule
\end{tabular}
\end{center}

Die Umgebungsvariablen werden beim Start über
\texttt{setup\_ros\_env.bat} gesetzt, das in
\texttt{isaac-sim-dobot.bat} (Abschnitt~\ref{sec:startskript})
eingebunden ist.

\subsection{Pythonpakete der Projekt-Entwicklungsumgebung}

Das Projekt-\texttt{.venv} (Python 3.14.3,
\texttt{dobot\_ros/.venv}) dient \textbf{ausschließlich}
für Entwicklungswerkzeuge (z.\,B. LaTeX-Build-Skript) und
enthält nur:

\begin{center}
\begin{tabular}{ll}
  \toprule
  Paket & Version \\
  \midrule
  \texttt{paho-mqtt} & 2.1.0 \\
  \texttt{pip} & 26.1 \\
  \bottomrule
\end{tabular}
\end{center}

Die eigentlichen Laufzeitabhängigkeiten (\texttt{pydobot},
\texttt{pyserial}, \texttt{rclpy}, \texttt{omni.*}) werden
\textbf{innerhalb der Isaac-Sim-Python-Umgebung} installiert bzw.
sind dort bereits enthalten.

\subsection{Laufzeitabhängigkeiten (Isaac-Sim-Umgebung)}

\begin{center}
\begin{tabular}{lll}
  \toprule
  Paket & Version & Bereitstellung \\
  \midrule
  \texttt{pydobot} & $\geq$ 1.3 & via \texttt{omni.kit.pipapi} (Install-Button) \\
  \texttt{pyserial} & $\geq$ 3.5 & via \texttt{omni.kit.pipapi} (Abhängigkeit) \\
  \texttt{rclpy} & Humble-Build & eingebettet in Isaac-Sim-Bridge \\
  \texttt{sensor\_msgs} & Humble-Build & eingebettet in Isaac-Sim-Bridge \\
  \texttt{omni.ui} & Kit-intern & bereits in Isaac Sim vorhanden \\
  \texttt{omni.ext} & Kit-intern & bereits in Isaac Sim vorhanden \\
  \texttt{omni.isaac.dynamic\_control} & Kit-intern & bereits in Isaac Sim vorhanden \\
  \texttt{omni.kit.pipapi} & Kit-intern & bereits in Isaac Sim vorhanden \\
  \bottomrule
\end{tabular}
\end{center}

\subsection{Systemanforderungen (Zusammenfassung)}

\begin{center}
\begin{tabular}{ll}
  \toprule
  Komponente & Mindestanforderung / empfohlene Version \\
  \midrule
  NVIDIA Isaac Sim & 5.1.0 oder neuer \\
  Isaac Sim Python & 3.11 (intern, kein manueller Eingriff) \\
  NVIDIA Treiber (Windows) & $\geq$ 537.58, empfohlen: 592.27+ \\
  CUDA Toolkit & 12.1 oder neuer (für GPU-Physik) \\
  GPU & NVIDIA RTX-fähige GPU, mind. 8\,GB VRAM \\
  RAM & mind. 16\,GB (empfohlen 32\,GB) \\
  Betriebssystem & Windows 11 (64-Bit) \\
  ROS2 & Humble oder Jazzy (embedded, kein Standalone-Install) \\
  Dobot Magician & USB-Verbindung, bekannter COM-Port \\
  \bottomrule
\end{tabular}
\end{center}

% ===================================================================
\section{Datei- und Paketstruktur}
\label{sec:dateien}
% ===================================================================

\subsection{Quellcode-Dateien}

\begin{center}
\begin{tabular}{lp{9.5cm}}
  \toprule
  Datei & Beschreibung \\
  \midrule
  \texttt{extension.py}       & Hauptlogik (1\,536 Zeilen): ROS-Node, Kinematik,
                                DC-Kopplung, UI, Teach-in, Würfel-Tracking \\
  \texttt{constants.py}       & Zentrale Farbwerte, USD-Pfad des Roboters,
                                Nullposition, Servo-Konfiguration \\
  \texttt{\_\_init\_\_.py}    & Exportiert \texttt{DobotBridgeExtension}
                                als Extension-Einstiegspunkt \\
  \texttt{extension.toml}     & Isaac-Sim-Manifest: Titel, Version, Python-Modul \\
  \texttt{install\_extension.py} & Hilfsskript: Kopiert Extension in den
                                Omniverse-Extensions-Ordner \\
  \texttt{set\_extension\_path.ps1} & PowerShell-Skript: Setzt
                                \texttt{OMNI\_KIT\_EXTENSION\_PATH} \\
  \texttt{build\_docs.py}     & Erstellt \texttt{documentation.pdf} aus LaTeX \\
  \texttt{build\_docs.bat}    & Windows-Batch-Wrapper für \texttt{build\_docs.py} \\
  \texttt{documentation.tex}  & Diese Dokumentation (LaTeX-Quelle) \\
  \bottomrule
\end{tabular}
\end{center}

\subsection{Extension-Manifest (\texttt{extension.toml})}

Die Datei \texttt{extension.toml} ist die \textbf{Visitenkarte} der Extension:
Sie teilt dem Isaac-Sim-Extension-Manager mit, wie die Extension heißt,
welche Version sie hat und welches Python-Modul als Einstiegspunkt dient.

\begin{warnbox}
\textbf{Häufigste Fehlerquelle beim Laden der Extension:}
Wird eine Extension im Extension-Manager nicht gefunden oder lässt sich nicht
korrekt laden, liegt das Problem fast immer in \texttt{extension.toml} ---
entweder stimmt der Dateiinhalt (Modulname, Pfad) nicht, oder der
Ordner, der \texttt{extension.toml} enthält, ist im Extension-Manager
nicht als Suchpfad eingetragen.
\end{warnbox}

\begin{lstlisting}[style=batch, caption={extension.toml}]
[package]
title       = "Dobot Digital Twin"
description = "ROS 2 Treiber und UI fuer den Dobot"
version     = "1.0.0"
authors     = ["Tobi"]

[dependencies]
"omni.kit.uiapp" = {}
"omni.ui"        = {}

[[python.module]]
name = "dobot_ros"
path = "."
\end{lstlisting}

Das Modul \texttt{dobot\_ros} wird aus dem Extension-Stammverzeichnis
geladen; \texttt{\_\_init\_\_.py} exportiert \texttt{DobotBridgeExtension},
die von Isaac Sim als \texttt{omni.ext.IExt}-Implementierung erkannt wird.

\subsection{Installations-Hilfsskript (\texttt{install\_extension.py})}

\begin{lstlisting}[caption={install\_extension.py -- Verwendung}]
# Standard-Installation in Omniverse-Roaming-Verzeichnis:
python install_extension.py

# Benutzerdefiniertes Zielverzeichnis:
python install_extension.py --dest "C:\MeinPfad" --name dobot_ros

# Vorschau ohne Dateikopie:
python install_extension.py --dry-run
\end{lstlisting}

Standardmäßig wird die Extension nach
\texttt{\%APPDATA\%\textbackslash{}NVIDIA Corporation\textbackslash{}Omniverse\textbackslash{}Extensions}
kopiert. Aus dem Quellbaum werden ausgeschlossen:
\texttt{.venv}, \texttt{.git}, \texttt{\_\_pycache\_\_},
\texttt{*.pyc/pyo/pyd}, kompilierte LaTeX-Artefakte.

\subsection{USD- und Robotermodell-Dateien}

\begin{center}
\begin{tabular}{lp{9.5cm}}
  \toprule
  Datei & Beschreibung \\
  \midrule
  \texttt{load\_dobot.urdf} & URDF-Roboterbeschreibung (162 Zeilen):
                              Gelenke, Massen, Kollisionsgeometrie \\
  \texttt{load\_dobot/load\_dobot.usd} & USD-Wurzeldatei,
                              führt alle Konfigurationsschichten zusammen \\
  \texttt{.../load\_dobot\_base.usd} & Mesh- und Kollisionsgeometrie
                              aller Roboterglieder (4,2\,MB) \\
  \texttt{.../load\_dobot\_robot.usd} & Artikulationsroot,
                              Gelenkeigenschaften, IsaacRobotAPI-Schema \\
  \texttt{.../load\_dobot\_physics.usd} & Physikparameter: Massen,
                              Trägheitstensoren, Dämpfung, Steifigkeit \\
  \texttt{.../load\_dobot\_sensor.usd} & Sensor-Platzhalter
                              (Erweiterungspunkt für zukünftige Sensoren) \\
  \texttt{meshes/*.STL, *.dae} & 3D-Netz-Dateien für alle Roboterglieder \\
  \bottomrule
\end{tabular}
\end{center}

% ===================================================================
\section{Systemarchitektur}
\label{sec:architektur}
% ===================================================================

\subsection{Schichtenmodell}

Die Extension ist in vier klar getrennte Schichten gegliedert:

\begin{center}
\renewcommand{\arraystretch}{1.4}
\begin{tabular}{|c|l|l|}
  \hline
  \textbf{Schicht} & \textbf{Klassen / Komponenten} & \textbf{Technologie} \\
  \hline
  UI-Schicht & \texttt{DobotBridgeExtension}: Fenster, Tabs,
               Callbacks & \texttt{omni.ui} \\
  \hline
  Simulations-Schicht & Dynamic-Control-API, USD-Stage,
                        omni.timeline & \texttt{omni.isaac.*} \\
  \hline
  Middleware-Schicht & ROS2-Node, Publisher,
                       \texttt{/joint\_states} & \texttt{rclpy} (Humble) \\
  \hline
  Hardware-Schicht & \texttt{DobotROSNode}: Pose-Abfrage,
                     Bewegungsbefehle & \texttt{pydobot}, pyserial \\
  \hline
\end{tabular}
\end{center}

\subsection{Hauptklassen}

\begin{description}[leftmargin=4.5cm, style=nextline]
  \item[\texttt{DobotROSNode}]
    Kapselt die serielle USB-Verbindung zum echten Dobot-Roboter
    über \texttt{pydobot} sowie den ROS2-Node und den
    JointState-Publisher. Enthält alle Hardware-nahen Methoden:
    Pose-Abfrage, Kinematikberechnung, Bewegungsbefehle,
    Vakuumgreifer, RC-Servo-Ansteuerung und
    Teach-in-Sequenzsteuerung.

  \item[\texttt{DobotBridgeExtension}]
    Subklasse von \texttt{omni.ext.IExt}. Verwaltet das
    Omniverse-UI-Fenster mit zwei Tabs, die Isaac-Sim-Dynamic-%
    Control-Kopplung, den Frame-Update-Loop sowie die
    Hintergrundthreads für Teach-in-Wiedergabe und Würfel-Tracking.
\end{description}

\subsection{Datenfluss}

\subsubsection{Richtung 1: Echter Roboter $\to$ Simulation (Digital Twin)}

\begin{enumerate}
  \item Pro Frame ruft \texttt{DobotBridgeExtension.\_on\_update()}
        die Methode \texttt{DobotROSNode.update\_step()} auf.
  \item \texttt{update\_step()} liest via
        \texttt{device.pose()} die aktuelle Pose:
        $(x, y, z, r, j_1, j_2, j_3, j_4)$ in mm und Grad.
  \item Kinematische Transformation
        (Abschnitt~\ref{sec:kinematik}) liefert
        Isaac-Sim-kompatible Gelenkwinkel.
  \item Konvertierung nach Radiant; Publikation als
        \texttt{JointState} auf \texttt{/joint\_states}.
  \item \texttt{\_drive\_joints()} setzt Winkel über
        DC-API auf die Isaac-Sim-Artikulation.
  \item Status-Dashboard im UI aktualisieren.
\end{enumerate}

\subsubsection{Richtung 2: Simulation / UI $\to$ Echter Roboter}

\begin{enumerate}
  \item Nutzer betätigt eine Steuertaste
        (z.\,B. \textit{X+ 5\,mm}).
  \item UI-Callback wird ausgelöst
        (z.\,B. \texttt{\_on\_jog\_clicked(5.0, 0.0)}).
  \item \texttt{DobotROSNode.jog(dx, dy)} wird aufgerufen.
  \item Aktuelle Pose einlesen, Delta addieren,
        \texttt{move\_to()} an Dobot senden.
  \item Nächster Frame liest neue Position -- Regelkreis geschlossen.
\end{enumerate}

% ===================================================================
\section{Lazy-Import-System}
\label{sec:lazyimport}
% ===================================================================

ROS2 (\texttt{rclpy}) und \texttt{pydobot} werden nicht beim
Extension-Start importiert, sondern erst beim ersten Klick auf
\textit{Connect}. Dieser Mechanismus verfolgt zwei Ziele:

\begin{enumerate}
  \item Die Extension lädt fehlerfrei, auch wenn ROS2 oder
        \texttt{pydobot} noch nicht installiert sind.
  \item Teure Import-Operationen verzögern den Isaac-Sim-Start nicht.
\end{enumerate}

Einmal erfolgreich geladene Module werden in Modul-globalen
Variablen gecacht:

\begin{lstlisting}[caption={Lazy-Import-Caches und \_try\_import\_ros()}]
_rclpy      = None   # rclpy-Modul
_Node       = None   # rclpy.node.Node
_JointState = None   # sensor_msgs.msg.JointState
_Dobot      = None   # pydobot.Dobot

def _try_import_ros():
    global _rclpy, _Node, _JointState
    if _rclpy is not None:
        return True        # Bereits gecacht
    try:
        import rclpy
        from rclpy.node import Node
        from sensor_msgs.msg import JointState
    except Exception:
        return False       # ROS2-Bridge nicht verfuegbar
    _rclpy, _Node, _JointState = rclpy, Node, JointState
    return True

def _try_import_dobot():
    global _Dobot
    if _Dobot is not None:
        return True
    try:
        from pydobot import Dobot
    except Exception:
        return False       # pydobot nicht installiert
    _Dobot = Dobot
    return True
\end{lstlisting}

% ===================================================================
\section{Kinematisches Modell}
\label{sec:kinematik}
% ===================================================================

\subsection{Koordinatenkonventionen (pydobot)}

Die Bibliothek \texttt{pydobot} liest vom Dobot Magician über das
serielle Protokoll (Kommando-ID\,10) die aktuelle Pose:
$(x, y, z, r, j_1, j_2, j_3, j_4)$.
Dabei gilt folgende Konvention:

\begin{align}
  j_1 &: \text{Basis-Rotation (Yaw, Z-Achse)} \quad [\text{Grad}] \\
  j_2 &: \text{Oberarm-Pitch, absoluter Weltwinkel von Horizontal}
          \quad [\text{Grad}] \\
  j_3 &: \text{Unterarm-Pitch, absoluter Weltwinkel von Horizontal}
          \quad [\text{Grad}] \\
  j_4 &: \text{Servo-gesteuerte Endeffektordrehung (Wristgelenk)}
          \quad [\text{Grad}]
\end{align}

\begin{notebox}
Durch den Parallelogramm-Mechanismus des Dobot Magician gilt
näherungsweise $j_3 \approx 0{{}^\circ}$ (Unterarm bleibt in
Weltkoordinaten horizontal), sofern das Seilzug- und Zahnradsystem
korrekt eingestellt ist.
\end{notebox}

\subsection{Isaac-Sim-Gelenkstruktur}

Das USD-Modell unter \texttt{/World/magician/base\_link} definiert
vier aktive Revolute-Gelenke:

\begin{center}
\begin{tabular}{clll}
  \toprule
  DOF-Index & USD-Name & Kinematische Bedeutung & Physischer Bereich \\
  \midrule
  0 & \texttt{joint\_1} & Basis-Yaw & $\pm 180{{}^\circ}$ \\
  1 & \texttt{joint\_2} & Schulter-Pitch & $0{{}^\circ}$--$90{{}^\circ}$ \\
  2 & \texttt{joint\_5} & Parallelogramm-Kompensation & $0{{}^\circ}$--$120{{}^\circ}$ \\
  3 & \texttt{joint\_6} & Endeffektor-Vertikalisierung & $\pm 180{{}^\circ}$ \\
  \bottomrule
\end{tabular}
\end{center}

\subsection{Einheitenstrategie}

Alle internen Berechnungen verwenden \textbf{Grad}.
Radiant-Konvertierung findet ausschließlich an den API-Grenzen statt:

\begin{center}
\begin{tabular}{lll}
  \toprule
  Kontext & Einheit & Konvertierung \\
  \midrule
  Interne Kinematik & Grad & -- \\
  DC-API \texttt{set\_dof\_state.pos} & Radiant & \texttt{math.radians()} \\
  DC-API \texttt{set\_dof\_position\_target} & Radiant & \texttt{math.radians()} \\
  ROS2 \texttt{JointState.position} & Radiant & \texttt{math.radians()} \\
  UI-Dashboard-Anzeige & Grad & -- \\
  \bottomrule
\end{tabular}
\end{center}

\subsection{Berechnungsformeln}
\label{sec:formeln}

\subsubsection{joint\_1 -- Gierwinkel mit festem Offset}

Der Gierwinkel $j_1$ des Dobot wird mit einem hardcodierten Offset von
$\Delta_\text{Yaw} = -45{{}^\circ}$ korrigiert, um die Ausrichtung der
Simulationsachse gegenüber dem realen Roboter auszugleichen:

\begin{equation}
  \theta_1 = j_1 + \Delta_\text{Yaw} = j_1 - 45{{}^\circ}
  \label{eq:j1}
\end{equation}

Der Wert ist als Modul-Konstante \texttt{J1\_YAW\_OFFSET\_DEG = -45.0}
in \texttt{extension.py} festgelegt.

\subsubsection{joint\_2 -- Direkte Übernahme}

\begin{align}
  \theta_2 &= j_2 \quad [{}^\circ]
\end{align}

\subsubsection{joint\_5 -- Parallelogramm-Kompensation}
\label{sec:j5}

Der Parallelogramm-Mechanismus hält \texttt{link\_5} in
Weltkoordinaten horizontal. $j_3$ und $j_2$ sind beide absolute
Weltwinkel; \texttt{joint\_5} muss ihre Differenz kodieren:

\begin{equation}
  \theta_5 = j_3 - j_2 \quad [{}^\circ]
  \label{eq:j5}
\end{equation}

\begin{lstlisting}[caption={Berechnung joint\_5 in update\_step()}]
# Parallelogramm-Kompensation: Unterarm in Weltlage halten
pos_j5 = j3 - j2   # Grad
\end{lstlisting}

\subsubsection{joint\_6 -- Endeffekteur-Vertikalisierung}
\label{sec:j6}

\paragraph{Anforderung:}
\texttt{link\_6} (Saugnapf-Träger) soll stets senkrecht nach unten
zeigen, unabhängig von der Armkonfiguration.

\paragraph{Herleitung:}
Die Summe aller Pitch-Rotationen entlang der kinematischen Kette
muss in Weltkoordinaten $90{{}^\circ}$ ergeben:

\begin{equation}
  \theta_2 + \theta_5 + \theta_6 = 90{{}^\circ}
  \label{eq:kette}
\end{equation}

Einsetzen von Gleichung~\eqref{eq:j5}:
\begin{align}
  j_2 + (j_3 - j_2) + \theta_6 &= 90{{}^\circ} \notag \\
  j_3 + \theta_6                &= 90{{}^\circ} \notag
\end{align}

Berücksichtigung der Servo-Verdrehung $j_4$:
\begin{equation}
  \boxed{\theta_6 = 90{{}^\circ} - j_3 - j_4}
  \label{eq:j6}
\end{equation}

\begin{warnbox}
\textbf{Wichtiger Hinweis:} Der Term $j_2$ tritt in
Gleichung~\eqref{eq:j6} \emph{nicht} auf -- er kürzt sich heraus.
Zwei Entwicklungsfehler wurden behoben:
\begin{enumerate}
  \item \texttt{pos\_j6 = j2 + j3 - j4} ließ $\theta_6$ mit der
        Schulterposition mitlaufen → Endeffektor kippte bei
        angehobenem Arm nach oben.
  \item Eine Mischung von $\frac{\pi}{2}$ (Radiant) und Grad-Werten
        in einer einzigen Formel führte zu Einheitenfehlern.
\end{enumerate}
\end{warnbox}

\begin{lstlisting}[caption={Korrekte Berechnung joint\_6 in update\_step()}]
# Kette: j2 + (j3-j2) + j6 = 90  =>  j6 = 90 - j3 - j4
# j2 darf hier NICHT auftreten
pos_j6 = 90.0 - j3 - j4   # Grad
\end{lstlisting}

\subsection{Vollständige Kinematik in \texttt{update\_step()}}

\begin{lstlisting}[caption={Kinematik-Kern in update\_step() (Ausschnitt aus extension.py)}]
pose = self.device.pose()  # (x, y, z, r, j1, j2, j3, j4)
if pose is not None and len(pose) >= 8:
    x, y, z, r = pose[0], pose[1], pose[2], pose[3]
    j1, j2, j3, j4 = pose[4], pose[5], pose[6], pose[7]

    # Gierwinkel-Offset: Simulationsachse gegenueber Realroboter
    j1 = j1 + J1_YAW_OFFSET_DEG   # -45 Grad (hardcoded)

    # Parallelogramm: Unterarm in Weltlage halten
    pos_j5 = j3 - j2
    # Endeffector vertikal: j2+(j3-j2)+j6=90 => j6=90-j3-j4
    pos_j6 = 90.0 - j3 - j4

    joints_deg = [j1, j2, pos_j5, pos_j6]
    self.last_valid_positions = joints_deg
    self.last_pose = {'x': x, 'y': y, 'z': z, 'r': r, 'j1': j1}

    # ROS2-Konvention: Radiant
    msg.position = [math.radians(v) for v in joints_deg]
    self.publisher_.publish(msg)
\end{lstlisting}

% ===================================================================
\section{Klasse \texttt{DobotROSNode}}
\label{sec:rosnode}
% ===================================================================

\subsection{Initialisierung (\texttt{\_\_init\_\_})}

Der Konstruktor führt in dieser Reihenfolge aus:
\begin{enumerate}
  \item Lazy-Imports prüfen; bei Fehler \texttt{RuntimeError} mit
        beschreibender Meldung.
  \item ROS2-Kontext initialisieren (\texttt{rclpy.init()},
        idempotent durch \texttt{rclpy.ok()}-Guard).
  \item ROS2-Node \texttt{dobot\_extension\_bridge} erstellen.
  \item JointState-Publisher auf \texttt{/joint\_states} anlegen
        (Queue-Tiefe\,10).
  \item Serielle Verbindung via \texttt{pydobot.Dobot(port=com\_port)}.
  \item Fallback-Puffer
        \texttt{last\_valid\_positions = [0.0, 0.0, 0.0, 0.0]}.
\end{enumerate}

\subsection{Methoden-Referenz}

\begin{longtable}{lp{9.5cm}}
  \toprule
  Methode & Beschreibung \\
  \midrule
  \endhead
  \texttt{update\_step()} &
    Liest Pose vom Dobot, berechnet Gelenkwinkel, publiziert
    \texttt{JointState}, ruft \texttt{rclpy.spin\_once(timeout=0)}
    auf (nicht-blockierend). Gibt
    \texttt{(pose\_dict, joints\_deg\_list)} zurück. Bei Fehler:
    Fallback auf \texttt{last\_valid\_positions}. \\[4pt]
  \texttt{set\_suction(active)} &
    Schaltet Vakuumgreifer SW1 (Kommando ID\,62). Fallback-Liste
    über mehrere pydobot-Methodennamen
    (\texttt{suck}, \texttt{set\_end\_effector\_suction},
    \texttt{set\_vacuum\_gripper}). \\[4pt]
  \texttt{move\_vertical(delta\_mm)} &
    Liest aktuelle Pose, addiert \texttt{delta\_mm} auf Z,
    sendet \texttt{move\_to}. Positiv = aufwärts. \\[4pt]
  \texttt{jog(dx, dy)} &
    Horizontaler Versatz in der X-Y-Ebene. Liest aktuelle Pose,
    addiert $(dx, dy)$ in mm. \\[4pt]
  \texttt{go\_home()} &
    Fährt zu \texttt{HOME\_X/Y/Z/R} aus \texttt{constants.py}. \\[4pt]
  \texttt{set\_servo\_gp3(angle\_deg)} &
    RC-Servo an GP3-PWM1 via Dobot-Rohprotokoll
    (Kommandos ID\,130 + ID\,132). Winkel $[0{{}^\circ},180{{}^\circ}]$,
    linearer Duty-Cycle. Details: Abschnitt~\ref{sec:servo}. \\[4pt]
  \texttt{play\_sequence(poses)} &
    Einfache Sequenz-Wiedergabe ohne Sicherheits-Anfahrt
    (interne Fallback-Methode; im Betrieb durch
    \texttt{\_play\_sequence\_thread} übernommen). \\[4pt]
  \texttt{stop\_sequence()} &
    Leert Dobot-interne Queue via
    \texttt{\_set\_queued\_cmd\_stop\_exec},
    \texttt{\_set\_queued\_cmd\_clear},
    \texttt{\_set\_queued\_cmd\_start\_exec}. \\[4pt]
  \texttt{shutdown()} &
    Schließt COM-Port via \texttt{device.close()},
    zerstört ROS2-Node via \texttt{destroy\_node()}. \\
  \bottomrule
\end{longtable}

\subsection{Fehlerbehandlung und Robustheit}

\paragraph{Kommunikationsunterbrechung:}
\texttt{update\_step()} fängt alle Ausnahmen auf. Bei Fehler werden
die Werte aus \texttt{last\_valid\_positions} verwendet -- das
Simulationsmodell \emph{bleibt in der letzten bekannten Stellung}.

\paragraph{pydobot-Versionsinkompatibilität:}
\begin{lstlisting}[caption={Suction-Fallback-Mechanismus in set\_suction()}]
for name, args in [
    ('suck',                     (active,)),
    ('set_end_effector_suction', (active,)),
    ('set_vacuum_gripper',       (1 if active else 0,)),
]:
    fn = getattr(self.device, name, None)
    if callable(fn):
        return fn(*args)
raise RuntimeError('Keine Saugbefehl-Methode in pydobot gefunden.')
\end{lstlisting}

% ===================================================================
\section{Servo-GP3-Ansteuerung}
\label{sec:servo}
% ===================================================================

\subsection{Hardware-Hintergrund}

Der Dobot Magician bietet einen GP3-Port (General Purpose Port 3)
mit einem PWM-fähigen Steuerausgang \texttt{PWM1} (IO-Adresse 15,
alternativ PWM2 = Adresse 16). Über diesen lässt sich ein Standard-RC-Servo
(Trägerfrequenz 50\,Hz, Pulsbreite 1--2\,ms) ansteuern.

\subsection{Duty-Cycle-Berechnung}

Der Duty-Cycle wird linear zwischen
$D_{\min} = 0{,}05$ (1\,ms-Puls, entspricht 0°) und
$D_{\max} = 0{,}10$ (2\,ms-Puls, entspricht 180°) interpoliert:

\begin{equation}
  D(\alpha) = D_{\min} + \frac{\alpha}{180{{}^\circ}} \cdot
              (D_{\max} - D_{\min}), \quad
  \alpha \in [0{{}^\circ},\, 180{{}^\circ}]
  \label{eq:duty}
\end{equation}

\subsection{Protokoll-Implementierung}

Da \texttt{pydobot} keine High-Level-API für GP3-PWM bereitstellt,
werden zwei Dobot-Rohprotokoll-Kommandos direkt gesendet:

\begin{lstlisting}[caption={set\_servo\_gp3() -- Rohprotokoll}]
angle_deg = max(0.0, min(180.0, float(angle_deg)))
duty = SERVO_MIN_DUTY + (angle_deg / 180.0) * \
       (SERVO_MAX_DUTY - SERVO_MIN_DUTY)

# Kommando 130: Pin 15 auf PWM-Modus (Multiplex-Typ 2)
mux = _Msg()
mux.id   = 130
mux.ctrl = 0x03    # Schreiben + in Queue
mux.params = bytearray()
mux.params.extend(struct.pack('B', GP3_PWM_ADDRESS))  # 15
mux.params.extend(struct.pack('B', 2))                # Typ: PWM
self.device._send_command(mux)

# Kommando 132: Frequenz und Duty-Cycle setzen
pwm = _Msg()
pwm.id   = 132
pwm.ctrl = 0x03
pwm.params = bytearray()
pwm.params.extend(struct.pack('B', GP3_PWM_ADDRESS))
pwm.params.extend(struct.pack('f', SERVO_FREQ_HZ))   # 50.0 Hz
pwm.params.extend(struct.pack('f', duty))
self.device._send_command(pwm)
\end{lstlisting}

\subsection{J1-Weltwinkel-Kompensation}
\label{sec:j1kompensation}

Beim Setzen des Servo-Winkels wird der \emph{Weltwinkel} gespeichert:
\begin{equation}
  \alpha_{\text{Welt}} = \alpha_{\text{Servo}} + j_1
\end{equation}

Bei jeder Fahrt in der Teach-in-Wiedergabe (\texttt{\_servo\_j1\_compensate()})
wird der Servo-Winkel rückgerechnet:
\begin{equation}
  \alpha_{\text{Servo,neu}} = \max\bigl(0{{}^\circ},\,
    \min(180{{}^\circ},\, \alpha_{\text{Welt}} - j_{1,\text{aktuell}})\bigr)
\end{equation}

Damit bleibt die Raumorientierung eines gehaltenen Objekts bei
wechselnder Basisrotation konstant.

% ===================================================================
\section{Dynamic-Control-Kopplung}
\label{sec:dc}
% ===================================================================

\subsection{Initialisierungsproblem}

Die Isaac-Sim-Dynamic-Control-API benötigt eine \emph{laufende}
Physik-Simulation. Vor dem Drücken von \emph{Play} liefert
\texttt{get\_articulation()} \texttt{INVALID\_HANDLE}.

\subsection{Lazy-Initialisierung}

\begin{lstlisting}[caption={Lazy DC-Init in \_on\_update()}]
if self._art_init_pending and self._dc is None:
    import omni.timeline
    if omni.timeline.get_timeline_interface().is_playing():
        robot_path = self._robot_path_input.model \
                         .get_value_as_string().strip()
        self._init_articulation(robot_path, silent=True)
        if self._dc is not None:
            self._art_init_pending = False
\end{lstlisting}

\subsection{DOF-Handles abrufen}

\begin{lstlisting}[caption={DOF-Handle-Abruf in \_init\_articulation()}]
joint_names = ['joint_1', 'joint_2', 'joint_5', 'joint_6']
for jname in joint_names:
    dof = dc.find_articulation_dof(art, jname)
    self._dof_handles.append(
        dof if dof != dc_mod.INVALID_HANDLE else None
    )
found = sum(1 for d in self._dof_handles if d is not None)
# Status: "Modell verbunden: 4/4 DOFs gefunden."
\end{lstlisting}

\subsection{Gelenkwinkel setzen -- \texttt{\_drive\_joints()}}

Pro Gelenk werden zwei komplementäre DC-Aufrufe verwendet:

\begin{lstlisting}[caption={\_drive\_joints() -- doppelter DC-Aufruf}]
for dof_handle, angle_deg in zip(self._dof_handles, joints_deg):
    if dof_handle is None:
        continue
    rad = math.radians(float(angle_deg))  # Grad -> Radiant

    # 1) Kinematischen Zustand direkt setzen
    state = self._dc_mod.DofState()
    state.pos = rad
    state.vel = 0.0
    self._dc.set_dof_state(dof_handle, state, self._dc_mod.STATE_POS)

    # 2) PhysX-Drive-Ziel setzen (Pflicht: [-2pi, 2pi])
    self._dc.set_dof_position_target(dof_handle, rad)
\end{lstlisting}

\begin{warnbox}
\texttt{set\_dof\_position\_target()} erzwingt $[-2\pi,\, 2\pi]$:
\texttt{PxArticulationJoint::setDriveTarget() only supports target
angle in range [-2Pi, 2Pi]}. Die Kinematik-Formeln halten alle
Ausgaben innerhalb dieses Bereichs.
\end{warnbox}

\subsection{Artikulations-USD-Pfad}

Die \texttt{ArticulationRootAPI} liegt beim importierten Dobot-Modell
auf \texttt{/World/magician/base\_link} -- \emph{nicht} auf dem
Elternprim \texttt{/World/magician}. Der vorausgefüllte Standardwert
im UI ist entsprechend gesetzt.

% ===================================================================
\section{Teach-in-System (Pick \& Place)}
\label{sec:teachin}
% ===================================================================

\subsection{Datenstruktur der Wegpunkte}

Jeder Wegpunkt ist ein 5-Tupel:
$(x_i,\, y_i,\, z_i,\, r_i,\, s_i) \in \mathbb{R}^4 \times \{0, 1\}$
mit $(x,y,z)$ in mm, $r$ in Grad, $s$ = Saugzustand (0 = aus, 1 = ein).

\subsection{Aufzeichnung}

\begin{description}[leftmargin=4.5cm, style=nextline]
  \item[\textit{Rec Punkt}]
    Speichert \texttt{DobotROSNode.last\_pose} + aktuellen
    virtuellen Saugzustand als neuen Wegpunkt.
    \emph{Kein Hardware-Befehl} -- der Saugzustand wird erst
    bei der Wiedergabe ausgeführt.

  \item[\textit{Saugen EIN/AUS}]
    Wechselt den virtuellen Saugzustand \emph{und} speichert
    automatisch einen Wegpunkt mit dem neuen Zustand
    (\texttt{\_on\_rec\_suction\_toggle\_clicked()}). Ermöglicht
    präzises Einteachen von Greif- und Ablegepositionen.

  \item[\textit{Löschen}]
    Setzt Wegpunktliste, Saugzustand und Zähler zurück.
\end{description}

\subsection{Wiedergabe-Thread}
\label{sec:playthread}

Die Wiedergabe läuft in einem dedizierten
\texttt{threading.Thread} (\texttt{daemon=True},
Name: \texttt{dobot\_play}).

\paragraph{Begründung für Threading statt asyncio:}
\texttt{pydobot.wait\_for\_cmd()} ist eine blockierende
Polling-Schleife. Im asyncio-Event-Loop würde sie den gesamten
Isaac-Sim-Hauptthread einfrieren und alle Frame-Updates blockieren.
Im Hintergrundthread blockiert sie nur diesen Thread;
Isaac Sim läuft ungehindert weiter. pydobots internes
\texttt{RLock} garantiert Thread-Sicherheit der seriellen Zugriffe.

\subsection{Bewegungsablauf pro Wegpunkt}

\begin{lstlisting}[caption={Kernlogik \_play\_sequence\_thread() (vereinfacht)}]
APPROACH_MM = 10.0   # 1 cm Sicherheitsabstand von oben

for i, (x, y, z, r, suction_active) in enumerate(poses):
    if self._stop_play_flag: break
    suction_changes = (suction_active != prev_suction)

    if suction_changes:
        # 1 cm ueber Zielpunkt anfahren (kein seitliches Streifen)
        _move_and_wait(x, y, z + APPROACH_MM, r)
        _servo_j1_compensate()
        _time.sleep(0.1)

    # Eigentliche Zielposition anfahren
    _move_and_wait(x, y, z, r)
    _servo_j1_compensate()

    # Saugzustand schalten (nach Positionserreichen)
    self._ros_node.set_suction(suction_active)
    _time.sleep(0.5)   # Druckaufbau/-abbau abwarten

    if suction_changes:
        # 1 cm sicher abheben
        _move_and_wait(x, y, z + APPROACH_MM, r)
        _servo_j1_compensate()

    prev_suction = suction_active
\end{lstlisting}

\subsection{Rückkehr zur Nullposition}

Nach dem letzten Wegpunkt oder nach Abbruch:
\begin{enumerate}
  \item 10\,mm über die Home-Position anfahren.
  \item Zur eigentlichen Home-Position absenken.
  \item Servo-Weltwinkel-Kompensation ausführen.
  \item Sauger deaktivieren.
\end{enumerate}

\subsection{Notstopp}

\texttt{\_on\_stop\_sequence\_clicked()} setzt
\texttt{\_stop\_play\_flag = True} und ruft
\texttt{DobotROSNode.stop\_sequence()} auf, das die interne
Dobot-Bewegungs-Queue leert.

% ===================================================================
\section{Ziel-Würfel-Tracking}
\label{sec:tracking}
% ===================================================================

\subsection{Würfel erstellen (\texttt{\_spawn\_target\_cube()})}

\begin{enumerate}
  \item Alten Würfel entfernen (Idempotenz).
  \item Startposition aus \texttt{DobotROSNode.last\_pose}
        (mm $\to$ m: Division durch 1000).
  \item \texttt{UsdGeom.Cube}, Kantenlänge 2,5\,cm,
        unter \texttt{/World/DobotTargetCube}.
  \item Orangenes \texttt{UsdPreviewSurface}-Material
        (diffuse RGB $(1.0, 0.45, 0.0)$, Opacity 0,9)
        unter \texttt{/World/DobotTargetCubeMat}.
  \item Tracking-Thread starten.
\end{enumerate}

\subsection{Tracking-Schleife (\texttt{\_tracking\_loop\_thread()})}

\begin{lstlisting}[caption={Tracking-Thread-Kernlogik}]
while not self._stop_tracking_flag and self._ros_node:
    pos = self._get_cube_world_position()  # m (USD-Stage)
    if pos is not None:
        x = pos[0] * 1000.0    # m -> mm
        y = pos[1] * 1000.0
        z = pos[2] * 1000.0
        r = self._ros_node.last_pose.get('r', 0.0)
        move_fn(x, y, z, r)   # Bewegungsbefehl an Dobot
    _time.sleep(1.0 / max(0.1, self._tracking_rate_hz))
\end{lstlisting}

Die Weltposition wird via
\texttt{UsdGeom.Xformable.ComputeLocalToWorldTransform()} exakt
ausgelesen (inkl. Elterntransformationen).
Tracking-Frequenz: 0,5--10,0\,Hz, einstellbar per Slider.

\subsection{Würfel entfernen (\texttt{\_despawn\_target\_cube()})}

Stop-Flag setzen → Thread mit \texttt{join(timeout=1\,s)} abwarten →
USD-Prims aus Stage entfernen → UI zurücksetzen.

% ===================================================================
\section{Klasse \texttt{DobotBridgeExtension}}
\label{sec:extension}
% ===================================================================

\subsection{Zustandsvariablen}

\begin{center}
\begin{tabular}{lp{9.5cm}}
  \toprule
  Variable & Bedeutung \\
  \midrule
  \texttt{\_ros\_node} & Aktive \texttt{DobotROSNode}-Instanz (oder \texttt{None}) \\
  \texttt{\_update\_sub} & Update-Event-Subscription (automatisch gefeuert) \\
  \texttt{\_suction\_active} & Aktueller Sauger-Status \\
  \texttt{\_servo\_angle} & Aktueller Servo-Winkel in Grad (Startwert 90°) \\
  \texttt{\_servo\_world\_angle} & Servo-Weltwinkel für J1-Kompensation \\
  \texttt{\_recorded\_poses} & Liste aller aufgezeichneten Wegpunkte \\
  \texttt{\_rec\_suction\_state} & Virtueller Saugzustand (nur Aufzeichnung) \\
  \texttt{\_play\_thread} & Hintergrundthread der Sequenz-Wiedergabe \\
  \texttt{\_stop\_play\_flag} & Abbruchsignal für Wiedergabe-Thread \\
  \texttt{\_cube\_spawned} & Würfel aktuell in Stage \\
  \texttt{\_tracking\_thread} & Hintergrundthread für Würfel-Tracking \\
  \texttt{\_stop\_tracking\_flag} & Abbruchsignal für Tracking-Thread \\
  \texttt{\_tracking\_rate\_hz} & Tracking-Sendefrequenz (Standard: 2,0\,Hz) \\
  \texttt{\_dc} & DC-Interface nach erfolgreicher Initialisierung \\
  \texttt{\_dc\_mod} & DC-Modul-Referenz für \texttt{DofState}-Konstruktor \\
  \texttt{\_art\_handle} & Artikulations-Handle \\
  \texttt{\_dof\_handles} & Liste von DOF-Handles [\texttt{j1,j2,j5,j6}] \\
  \texttt{\_art\_init\_pending} & True nach Connect, False nach erfolgreicher DC-Init \\
  \texttt{\_dock\_task} & Async-Docking-Task \\
  \bottomrule
\end{tabular}
\end{center}

\subsection{Lebenszyklusmethoden}

\begin{description}[leftmargin=4.5cm, style=nextline]
  \item[\texttt{on\_startup(ext\_id)}]
    Alle Zustandsvariablen initialisieren; UI-Fenster
    (490$\times$590\,px) aufbauen; Docking-Task starten.

  \item[\texttt{\_dock\_window\_deferred()}]
    Async-Coroutine: 2 Frames warten, Fenster neben
    Property-Panel docken via \texttt{ui.DockPosition.SAME}.
    Zwei Frames notwendig, da Fenster erst nach erstem
    Render im Workspace-Index erscheint.

  \item[\texttt{\_on\_update(event)}]
    Frame-Callback: DC-Lazy-Init, Dobot-Polling,
    Gelenkwinkeltransfer, Dashboard-Update.

  \item[\texttt{on\_shutdown()}]
    Dock-Task canceln (verhindert ``extension object is still
    alive''-Warnung), \texttt{\_cleanup()} aufrufen,
    Fenster zerstören.
\end{description}

\subsection{Frame-Update-Ablauf (\texttt{\_on\_update()})}

\begin{lstlisting}[caption={\_on\_update() -- vollständiger Ablauf}]
def _on_update(self, event):
    if not self._ros_node:
        return

    # 1. Lazy DC-Initialisierung
    if self._art_init_pending and self._dc is None:
        import omni.timeline
        if omni.timeline.get_timeline_interface().is_playing():
            robot_path = self._robot_path_input.model \
                             .get_value_as_string().strip()
            self._init_articulation(robot_path, silent=True)
            if self._dc is not None:
                self._art_init_pending = False

    # 2. Dobot-Daten lesen und ROS2 publizieren
    pose, joints = self._ros_node.update_step()

    # 3. DC-Kopplung (nur wenn Simulation laeuft)
    import omni.timeline
    if omni.timeline.get_timeline_interface().is_playing():
        self._drive_joints(joints)

    # 4. Dashboard aktualisieren
    self._update_display(pose, joints)
\end{lstlisting}

\subsection{Ressourcenfreigabe (\texttt{\_cleanup()})}

\begin{enumerate}
  \item Stop-Flag Wiedergabe-Thread setzen;
        \texttt{join(timeout=1\,s)}.
  \item Stop-Flag Tracking-Thread setzen;
        \texttt{join(timeout=1\,s)}.
  \item Würfel-Prims aus USD-Stage entfernen.
  \item Update-Subscription freigeben.
  \item DC-Interface und Handles loslassen.
  \item \texttt{DobotROSNode.shutdown()} aufrufen
        (COM-Port schließen, ROS2-Node zerstören).
\end{enumerate}

% ===================================================================
\section{UI-Fenster}
\label{sec:ui}
% ===================================================================

\subsection{Tab 0: Steuerung}

\begin{center}
\begin{tabular}{lp{8cm}}
  \toprule
  UI-Element & Funktion \\
  \midrule
  COM-Port-Feld & Serieller Port (Standard: \texttt{COM3}) \\
  Connect-Button & Verbindungsaufbau; zeigt nach Erfolg grün
                   \texttt{$\bullet$ Verbunden} \\
  Disconnect-Button & Verbindungstrennung und UI-Reset \\
  Verbindungsstatus (\texttt{$\bullet$}) & Rot = getrennt, grün = verbunden \\
  Robot-USD-Feld & USD-Pfad der Artikulations-Root in der Szene;
                   Standard \texttt{/World/magician/base\_link} \\
  Install packages & Installiert \texttt{pydobot} und
                      \texttt{pyserial} via \texttt{omni.kit.pipapi} \\
  Suction ON/OFF & Vakuumgreifer umschalten \\
  Up / Down 10\,mm & Vertikalbewegung $\pm 10$\,mm \\
  X$\pm$ / Y$\pm$ 5\,mm & Horizontaler JOG $\pm 5$\,mm \\
  Servo GP3 $\pm 10{{}^\circ}$ & Servo-Winkel erhöhen / verringern;
                                   Winkel-Anzeige-Label \\
  Würfel spawnen / entfernen & Ziel-Würfel ein-/ausblenden \\
  Rate-Slider (0,5--10\,Hz) & Tracking-Frequenz konfigurieren \\
  Rate-Label & Zeigt aktuelle Tracking-Frequenz \\
  \bottomrule
\end{tabular}
\end{center}

\subsection{Tab 1: Pick \& Place}

\begin{center}
\begin{tabular}{lp{8cm}}
  \toprule
  UI-Element & Funktion \\
  \midrule
  Go Home & Roboter zur Nullposition fahren;
             Button-Label zeigt $x$/$y$/$z$-Koordinaten \\
  Servo GP3 $\pm 10{{}^\circ}$ & Identisch zu Tab\,0 \\
  Saugen (Aufn.) EIN/AUS & Virtueller Sauger-Toggle für Aufzeichnung;
                            schaltet Zustand um \emph{und} speichert
                            automatisch Wegpunkt \\
  Rec Punkt & Aktuellen Wegpunkt manuell aufzeichnen \\
  $n$\,Pkt. & Zähler der gespeicherten Wegpunkte \\
  Löschen & Alle Wegpunkte und Saugzustand zurücksetzen \\
  Play & Teach-in-Sequenz im Hintergrundthread starten \\
  Stop & Laufende Sequenz abbrechen und Dobot-Queue leeren \\
  Scrollbare Waypoint-Liste & Zeigt alle Punkte mit Koordinaten und
                               Saugzustand; Saugpunkte in Grün \\
  \bottomrule
\end{tabular}
\end{center}

\subsection{Status-Dashboard (immer sichtbar)}

\begin{center}
\begin{tabular}{lp{9cm}}
  \toprule
  Element & Beschreibung \\
  \midrule
  Status-Label & Farbkodierter Betriebsstatus:
                  \textcolor{green!50!black}{grün} = OK /
                  Verbunden,
                  \textcolor{orange!80!black}{gelb} = Pending /
                  Warnung,
                  \textcolor{red!70!black}{rot} = Fehler \\
  ROS-Topic-Label & \texttt{ROS topic: /joint\_states} \\
  Pose-Label & Kartesische Position:
                \texttt{Pose: x=... y=... z=... r=...} \\
  Joints-Label & Gelenkwinkel in Grad:
                  \texttt{Joints [°]: j1=... j2=... j5=... j6=...} \\
  \bottomrule
\end{tabular}
\end{center}

% ===================================================================
\section{Threading-Modell}
\label{sec:threading}
% ===================================================================

\begin{center}
\begin{tabular}{llp{7cm}}
  \toprule
  Thread & Name & Aufgabe \\
  \midrule
  Hauptthread & (Isaac Sim) & UI-Callbacks, Frame-Updates
                (\texttt{\_on\_update}), DC-API-Aufrufe \\
  Wiedergabe & \texttt{dobot\_play} & Blockierende
                \texttt{wait\_for\_cmd()}-Schleife \\
  Tracking & \texttt{dobot\_cube\_track} & Periodische
                \texttt{move\_to}-Befehle für Würfel \\
  \bottomrule
\end{tabular}
\end{center}

\paragraph{Thread-Sicherheit:}
pydobots internes \texttt{RLock} serialisiert alle seriellen
Zugriffe. Statusschreibzugriffe aus Hintergrundthreads
(\texttt{\_set\_status()}) sind durch den GIL (Global Interpreter
Lock) von CPython atomisch.

\paragraph{Graceful Shutdown:}
Beide Stop-Flags werden in \texttt{\_cleanup()} gesetzt;
danach \texttt{join(timeout=1\,s)} pro Thread.

% ===================================================================
\section{ROS2-Integration}
\label{sec:ros2}
% ===================================================================

\subsection{Node-Konfiguration}

\begin{center}
\begin{tabular}{ll}
  \toprule
  Parameter & Wert \\
  \midrule
  Node-Name & \texttt{dobot\_extension\_bridge} \\
  Topic & \texttt{/joint\_states} \\
  Nachrichtentyp & \texttt{sensor\_msgs/msg/JointState} \\
  Queue-Tiefe & 10 \\
  Publish-Frequenz & $\approx 60$\,Hz (Frame-Rate-gekoppelt) \\
  RMW & \texttt{rmw\_fastrtps\_cpp} (Fast DDS) \\
  ROS2-Distro & Humble (eingebettet in Isaac-Sim-Bridge) \\
  \bottomrule
\end{tabular}
\end{center}

\subsection{Nachrichtenstruktur}

\begin{lstlisting}[caption={JointState-Nachricht in update\_step()}]
msg.header.stamp = self._node.get_clock().now().to_msg()
msg.name     = ['joint_1', 'joint_2', 'joint_5', 'joint_6']
msg.position = [math.radians(j1),     # Radiant (ROS2-Konvention)
                math.radians(j2),
                math.radians(pos_j5),
                math.radians(pos_j6)]
# velocity und effort bleiben leer (optionale Felder)
\end{lstlisting}

\subsection{ROS2-Kontext-Verwaltung}

\texttt{rclpy.init()} wird genau einmal aufgerufen
(\texttt{rclpy.ok()}-Guard verhindert Mehrfach-Init).
\texttt{rclpy.spin\_once(timeout\_sec=0.0)} am Ende von
\texttt{update\_step()} verarbeitet eingehende Nachrichten
nicht-blockierend.

% ===================================================================
\section{Konstanten (\texttt{constants.py})}
\label{sec:konstanten}
% ===================================================================

\subsection{Farbwerte (ARGB-Hex)}

\begin{center}
\begin{tabular}{lll}
  \toprule
  Konstante & Hexwert & Verwendung \\
  \midrule
  \texttt{CLR\_BG\_DARK}  & \texttt{0xFF0A0F1A} & Sehr dunkler Hintergrund \\
  \texttt{CLR\_BG\_MID}   & \texttt{0xFF101828} & Aktives Panel / Toolbar \\
  \texttt{CLR\_ACCENT}    & \texttt{0xFF4A9EFF} & Info / Akzentfarbe \\
  \texttt{CLR\_GREEN}     & \texttt{0xFF4ADE80} & Erfolg / Verbunden \\
  \texttt{CLR\_RED}       & \texttt{0xFFEF6B6B} & Fehler / Getrennt \\
  \texttt{CLR\_YELLOW}    & \texttt{0xFFE0B040} & Warnung / Pending \\
  \texttt{CLR\_TEXT}      & \texttt{0xFFD0D8E8} & Primärtext \\
  \texttt{CLR\_TEXT\_DIM} & \texttt{0xFF607090} & Sekundärtext \\
  \texttt{CLR\_BORDER}    & \texttt{0xFF1A2840} & Trennlinien \\
  \bottomrule
\end{tabular}
\end{center}

\subsection{Nullposition und Servo}

\begin{center}
\begin{tabular}{lrl}
  \toprule
  Konstante & Wert & Beschreibung \\
  \midrule
  \texttt{HOME\_X} & $-3{,}7$\,mm & Nullposition X \\
  \texttt{HOME\_Y} & $227{,}8$\,mm & Nullposition Y \\
  \texttt{HOME\_Z} & $66{,}1$\,mm & Nullposition Z \\
  \texttt{HOME\_R} & $88{,}4{{}^\circ}$ & Nullposition Wristdrehung \\
  \texttt{GP3\_PWM\_ADDRESS} & 15 & IO-Adresse PWM1 am GP3-Port \\
  \texttt{SERVO\_FREQ\_HZ} & $50{,}0$\,Hz & RC-Servo-Trägerfrequenz \\
  \texttt{SERVO\_MIN\_DUTY} & $0{,}05$ & 1\,ms-Puls $= 0°$ \\
  \texttt{SERVO\_MAX\_DUTY} & $0{,}10$ & 2\,ms-Puls $= 180°$ \\
  \bottomrule
\end{tabular}
\end{center}

\subsection{USD-Pfade}

\begin{center}
\begin{tabular}{lp{8cm}}
  \toprule
  Konstante & Pfad \\
  \midrule
  \texttt{ROBOT\_USD\_PATH} & \texttt{/World/magician/base\_link} \\
  \bottomrule
\end{tabular}
\end{center}

% ===================================================================
\section{URDF-Robotermodell (\texttt{load\_dobot.urdf})}
\label{sec:urdf}
% ===================================================================

Das URDF (Unified Robot Description Format) beschreibt die
kinematische Struktur des Dobot Magician für den Isaac-Sim-Import.

\subsection{Gelenkstruktur}

\begin{center}
\begin{tabular}{llll}
  \toprule
  Gelenk & Typ & Achse & Limits \\
  \midrule
  \texttt{joint\_1} & revolute & Z (Yaw) &
    $\pm 180{{}^\circ}$, $\pm 10$\,rad/s \\
  \texttt{joint\_2} & revolute & Y (Schulter) &
    $0{{}^\circ}$--$90{{}^\circ}$, $\pm 10$\,rad/s \\
  \texttt{joint\_5} & revolute & Y (Parallelogramm) &
    $0{{}^\circ}$--$120{{}^\circ}$, $\pm 10$\,rad/s \\
  \texttt{joint\_6} & revolute & Y (Endeffektor) &
    $\pm 180{{}^\circ}$, $\pm 10$\,rad/s \\
  \texttt{joint\_suction} & fixed & -- & -- \\
  \bottomrule
\end{tabular}
\end{center}

\subsection{Körper (Links)}

\begin{center}
\begin{tabular}{ll}
  \toprule
  Link & Masse \\
  \midrule
  \texttt{base\_link} & 1{,}0\,kg \\
  \texttt{link\_1} -- \texttt{link\_4} & je 0{,}3\,kg \\
  \texttt{link\_5}, \texttt{link\_6} & 0{,}1\,kg \\
  \texttt{link\_suction} & 0{,}05\,kg \\
  \bottomrule
\end{tabular}
\end{center}

% ===================================================================
\section{Installation und Inbetriebnahme}
\label{sec:installation}
% ===================================================================

\subsection{Erstinstallation}

\begin{enumerate}
  \item \textbf{Isaac Sim starten:}
        \texttt{C:\textbackslash{}NVIDIA\_Isaac\_Sim\textbackslash{}isaac-sim-dobot.bat}
        aufrufen (setzt \texttt{OMNI\_KIT\_EXTENSION\_PATH}
        automatisch).

  \item \textbf{Extension aktivieren:}
        \textit{Window $\to$ Extensions}; nach
        \texttt{dobot\_ros} suchen und aktivieren.

  \item \textbf{Pakete installieren:}
        Im Extension-Panel \textit{Install packages} klicken.
        Dies installiert \texttt{pydobot} und \texttt{pyserial}
        in der Isaac-Sim-Python-Umgebung via
        \texttt{omni.kit.pipapi}.
        Danach Extension deaktivieren und neu aktivieren.
\end{enumerate}

\subsection{Betrieb Schritt für Schritt}

\begin{enumerate}
  \item Szene mit Dobot-Magician-Modell laden.
  \item \texttt{ArticulationRootAPI} auf \texttt{base\_link}
        prüfen (\textit{Property $\to$ Physics $\to$ Articulation Root}).
  \item COM-Port eingeben (Standard: \texttt{COM3}) und
        \textbf{Connect} klicken.
        Status: \textit{Verbunden -- warte auf Play\ldots}
  \item \textbf{Play} drücken.
        Status: \textit{Modell verbunden: 4/4 DOFs gefunden.}
  \item Simulationsmodell folgt echtem Roboter in Echtzeit.
\end{enumerate}

\subsection{Teach-in-Workflow}

\begin{enumerate}
  \item Tab \textit{Pick \& Place} wählen.
  \item Roboter manuell in Ausgangsposition fahren (Jog-Tasten).
  \item \textit{Rec Punkt} klicken.
  \item An Greifpunkt: \textit{Saugen (Aufn.)} auf EIN --
        automatisch Waypoint gespeichert.
  \item Schritte 2--4 für alle Punkte wiederholen.
  \item \textit{Play} starten; Sequenz läuft mit Sicherheits-Anfahrt.
  \item \textit{Stop} bricht jederzeit ab.
\end{enumerate}

% ===================================================================
\section{Bekannte Einschränkungen}
\label{sec:einschraenkungen}
% ===================================================================

\begin{itemize}
  \item \textbf{ROS2 eingebettet:}
        Die Extension verwendet die in Isaac Sim eingebettete
        ROS2 Humble Bridge. Eine separate ROS2-Installation ist
        nicht erforderlich, jedoch auch nicht kompatibel ohne
        Anpassung der Bibliothekspfade.

  \item \textbf{DC-API vor Play:}
        \texttt{get\_articulation()} vor dem ersten Play-Ereignis
        erzeugt eine Warnung. Durch Lazy-Init tritt diese nur
        einmal auf.

  \item \textbf{Suction-Methoden-Varianz:}
        Die Fallback-Liste deckt bekannte pydobot-Versionen ab;
        bei zukünftigen Versionen ggf. erweitern.

  \item \textbf{Servo-Zustand nach Reconnect:}
        Nach Disconnect/Reconnect startet der Servo-Winkel bei
        90° -- der tatsächliche Hardware-Winkel kann abweichen.

  \item \textbf{Keine Greifer-Rückmeldung:}
        Saugbefehle werden als Fire-and-Forget gesendet.
        Ob ein Objekt tatsächlich gegriffen ist, kann ohne
        Drucksensor nicht verifiziert werden.

  \item \textbf{Physikalisch ungültige Winkel:}
        Winkel außerhalb $[-2\pi,\, 2\pi]$ führen zum
        PhysX-Fehler in \texttt{set\_dof\_position\_target()}.
        Die Kinematik-Formeln halten alle Ausgaben im gültigen
        Bereich.
\end{itemize}

% ===================================================================
\section{Entwicklungsverlauf}
\label{sec:verlauf}
% ===================================================================

\begin{center}
\begin{tabular}{lp{9cm}}
  \toprule
  Commit & Meilenstein \\
  \midrule
  \texttt{1611913} & Initiales Repository-Setup \\
  \texttt{d2744a9} & ROS2-Bridge-Integration und Merge \\
  \texttt{1ed673f} & Schrittmotor-Kopplung funktionsfähig \\
  \texttt{d099a83} & Live/Sim-Umschaltung implementiert \\
  \texttt{c919311} & Vakuumgreifer-Ansteuerung funktionsfähig \\
  \texttt{7ddfb27} & Squeeze-/Pick-Logik funktionsfähig \\
  \texttt{a95763a} & Deckel-Maze-Bewegung implementiert \\
  \texttt{76b73a6} & MQTT-Integration für Deckel/Maze \\
  \texttt{72ec561} & MQTT für Digital-State, -Move und Binary-Move \\
  \texttt{1ef5eca} & Stabiler Arbeitsstand \\
  \texttt{7471a6f} & Vollständig kommentierter Digital-Twin \\
  \texttt{685be33} & Fehlerbehebungen (Kinematik-Formeln) \\
  \texttt{5cd2262} & Pick-\&-Place vollständig funktionsfähig (aktuell) \\
  \bottomrule
\end{tabular}
\end{center}

\end{document}
